% !TeX program = lualatex
% Compile twice: lualatex 50.tex
% All references and prose are embedded; no BibTeX database or external inputs.
\documentclass[11pt,a4paper]{article}
\usepackage[margin=24mm,headheight=14pt]{geometry}
\usepackage{fontspec}
\setmainfont{Latin Modern Roman}
\setsansfont{Latin Modern Sans}
\setmonofont{Latin Modern Mono}
\usepackage{amsmath,amssymb,mathtools}
\usepackage{unicode-math}
\setmathfont{Latin Modern Math}
\usepackage{microtype}
\usepackage{enumitem,xurl,needspace}
\usepackage[dvipsnames]{xcolor}
\usepackage{hyperref}
\hypersetup{unicode=true,colorlinks=true,linkcolor=MidnightBlue,urlcolor=MidnightBlue,
 pdftitle={Problem 50: One-Way Permutations -- Research Attempt},pdfauthor={Research notebook},bookmarksnumbered=true}
\usepackage{bookmark}
\usepackage{tocloft}
\setlength{\cftsecnumwidth}{3em}
\cftsetpnumwidth{2.5em}
\cftsetrmarg{4em}
\usepackage{titlesec}
\titleformat{\section}{\Large\bfseries\raggedright}{\thesection}{0.7em}{}
\usepackage{fancyhdr}
\pagestyle{fancy}\fancyhf{}
\fancyhead[L]{\small\sffamily Open Problems Math | Research notebook}
\fancyhead[R]{\small 27 September 2026}
\fancyfoot[C]{\thepage}
\setlength{\parindent}{0pt}
\setlength{\parskip}{5pt plus 1pt}
\setlength{\emergencystretch}{3em}
\setcounter{tocdepth}{1}
\setcounter{secnumdepth}{1}
\setlist[itemize]{leftmargin=3em,itemsep=5pt,topsep=4pt,parsep=0pt}
\allowdisplaybreaks
\newcommand{\N}{\mathbb{N}}
\newcommand{\Z}{\mathbb{Z}}
\newcommand{\Q}{\mathbb{Q}}
\newcommand{\R}{\mathbb{R}}
\newcommand{\C}{\mathbb{C}}
\newcommand{\F}{\mathbb{F}}
\newcommand{\A}{\mathbb{A}}
\newcommand{\PP}{\mathbb P}
\newcommand{\E}{\mathbb{E}}
\newcommand{\eps}{\varepsilon}
\newcommand{\dd}{\,\mathrm d}
\newcommand{\sref}[2]{\hyperlink{source:#1:#2}{[S#2]}}
\newcommand{\eref}[2]{\hyperlink{extra:#1:#2}{[E#2]}}
% Turn the next switch off for a shorter reading copy. The quotations preserve
% every exactTarget field, including qualifications omitted from a short summary.
\newif\ifshowcatalogue
\showcataloguetrue
\newcommand{\cataloguescope}[1]{\ifshowcatalogue
 \paragraph{Exact scope in the supplied catalogue.}
 \begin{quote}\small #1\end{quote}\fi}
\numberwithin{equation}{section}
\begin{document}
\setcounter{section}{49}
% Standalone export: only Problem 50 is present at or after 50 in the available source.
% BEGIN ORIGINAL SECTION 50
% ================================================================
% problemId: problem.existence-of-one-way-permutations
% source releaseRank: 50; source releaseStatus: open
% exactTarget SHA-256: fd2430866809dbadb2360110e66e50f8014362a28dfb0795b02d7c1309a9414d
\clearpage
\section{\texorpdfstring{Existence of One-Way Permutations}{Existence of One-Way Permutations}}\label{50}
\subsection{Definitions and mathematical statement}
For every $n\ge1$, require a bijection $f_n:\{0,1\}^n\to\{0,1\}^n$, with one polynomial-time algorithm evaluating $f_n(x)$ from $(1^n,x)$. For each probabilistic polynomial-time $A$, require
\[
 \Pr_{x\leftarrow\{0,1\}^n,A}[A(1^n,f_n(x))=x]\le\eta_A(n),
\]
where $\eta_A$ is negligible. Since $f_n$ is bijective, each output has exactly one preimage. This is an unconditional existence problem on all binary strings of each length, without sampled keys or variable domains.
\par\smallskip\textit{Formulation and concept references:} \sref{50}{1}.
\cataloguescope{There exists a family f\_n:\{0,1\}\textasciicircum{}n->\{0,1\}\textasciicircum{}n, n>=1, of bijections computed by one deterministic polynomial-time algorithm on (1\textasciicircum{}n,x), such that for every probabilistic polynomial-time algorithm A its inversion success Pr[A(1\textasciicircum{}n,f\_n(x))=x], over uniform x and A's coins, is negligible in n: for every c>0 it is at most n\textasciicircum{}(-c) for all sufficiently large n. This is unconditional existence in the classical unrelativized model, with no key sampling or variable domains.}
\subsection{Short English statement}
Can a length-preserving bijection be easy to compute but negligibly likely to be inverted by any efficient randomized algorithm?
% BEGIN RESEARCH INSERTION 50
\subsection{Research attempt}

\paragraph{Current frontier.}
\textit{Research date: 27 September 2026.}
The distinction between a single length-preserving permutation and a
collection indexed by sampled public keys is essential here. Goldreich's
clarification distinguishes full-cube index sets from merely samplable
index sets, and explicitly warns that the inverter need not receive the
coins used to generate an index. The present target is the single,
uniformly computable, full-cube version, not a keyed version with a hidden
setup tape. \eref{50}{1}

The inherited formulation reference \sref{50}{1} is a separation result:
Asharov--Segev rule out the specified fully black-box, domain-invariant
construction from one-way functions and indistinguishability obfuscation
for oracle-aided circuits. This is an established restriction on a
construction model, not a disproof of the present existence statement.
The publisher's abstract and the authors' presentation were consulted;
this notebook does not claim to have independently audited their full
separation proof. \eref{50}{2}

A relevant positive development is the \emph{conditional} full-input-domain
trapdoor-permutation construction in Shmueli--Zhandry's 2025 manuscript:
Theorem~6 assumes sub-exponentially secure indistinguishability obfuscation
and one-way functions. Their Section~1.4 distinguishes input-domain
invariance from key-domain invariance, explaining why its conclusion does
not contradict the preceding barrier. Construction~41 still generates a
public key and a secret key. Neither this hypothesis package nor its
sampled-key conclusion supplies the unconditional family required here.
The cited theorem is treated as a manuscript result, not independently
certified by this notebook. \sref{50}{2}, \eref{50}{3}

Two further facts delimit the attack. Direct-product amplification of weak
one-wayness is classical; the proof below supplies the quantitative step
needed here rather than presuming independence of an inverter's outputs.
\eref{50}{4} The 2025 in-place-query result gives a quantum permutation
inverter using $O(\sqrt{M})$ queries on an $M$-element domain. With
$M=2^n$, this is not a polynomial-in-$n$ algorithm, and a query theorem does
not establish classical, code-aware computational hardness.
\sref{50}{3}, \eref{50}{5} None of these inspected results supplies an
unconditional solution of the exact entry.

\paragraph{Attack route.}
Four routes were considered. First, compile a sampled-key construction
into one full-cube map; the missing ingredient is a way to obtain useful
public descriptions without also revealing their generation trapdoors.
Second, allow hardness only on a recognizable, inverse-polynomial-density
part of the cube, then combine a bijective extension with hardness
amplification; the missing ingredient becomes an explicit hard partial
action rather than a perfect full-domain construction. Third, obtain
hardness from cycle geometry or from elementary reversible gates; the
issue is whether the proposed structure already gives an inverse.
Fourth, prove a direct lower bound for the unique inverse's bits; this
would require an average-case lower bound inside unambiguous complexity
classes. The second route yields a fully specified conditional
normalization theorem. The other routes are pursued as construction
attempts and adversarial tests, not as additional catalogue problems.

\paragraph{Attempt.}
Throughout, $U_m$ is uniform on $\{0,1\}^m$. Algorithms are uniform and
have a polynomial worst-case running-time bound, including over their
random coins. A wrong-length output or a distinguished failure symbol
$\bot$ counts as failure. Quantities called negligible satisfy the
all-sufficiently-large-lengths convention of the original statement.
All lemmas below are proved in this notebook; historical novelty is not
asserted, and the amplification mechanism is explicitly classical.

\medskip\needspace{5\baselineskip}
\textbf{A. Attempting to remove the generated key.}
Write a deterministic execution of a randomized key generator as
\[
 (\operatorname{pk}_m(r),\operatorname{sk}_m(r))
       =\operatorname{Gen}(1^m;r),\qquad r\in\{0,1\}^{\ell_m}.
\]
Assume, even favorably, that every seed produces a valid permutation
$P_{m,r}$ of $\{0,1\}^m$, that $\ell_m$ is polynomially bounded, and that
both $P_{m,r}$ and its inverse can be evaluated in polynomial time given
$r$ (the inverse uses the regenerated secret key). These assumptions
only strengthen this candidate construction's correctness.
The input/output interface of Construction~41 motivates this test.
\eref{50}{3}

The obvious seed-preserving packing is
\begin{equation}\label{eq:50:seed}
 B_m(r,x)=(r,P_{m,r}(x)).
\end{equation}
It is a bijection, but its inverse is explicitly
\[
 B_m^{-1}(r,y)=(r,P_{m,r}^{-1}(y)).
\]
The inverter regenerates the secret key from $r$. Thus keeping the seed,
although it solves the encoding and bijectivity problems, destroys the
intended hardness. Conversely, discarding the seed is not a generally
valid injectivity argument; different seeds can generate the same public
description. Sampling a public key cannot simply be replaced by fixing a
convenient seed: the fixed seed is then reproducible by the inverter.
Existence of good keys also does not exhibit a polynomial-time procedure
selecting a good key at every length. These are separate issues from the
security of a correctly generated public key.

Could one \emph{hide} the seed instead? The following exact calculation
shows the circularity of a broad such proposal.

\textbf{Lemma A (masked-seed equivalence).}
Let $h_m:\{0,1\}^{\ell_m}\to\{0,1\}^{\ell_m}$ be efficiently evaluable and put
\begin{equation}\label{eq:50:masked}
 F_m(r,x)=\bigl(h_m(r),P_{m,r}(x)\bigr).
\end{equation}
Then $F_m$ is a bijection if and only if $h_m$ is a bijection. In that
case an inverter for $F_m$ gives an inverter for $h_m$ with at least the
same average success, while an inverter for $h_m$ gives an inverter for
$F_m$ with the same average success. The reductions have polynomial
running time in $m+\ell_m$ and use the same parameter $1^m$.

\textit{Proof.}
For any $(z,y)$ the number of its preimages under $F_m$ is exactly
$\#h_m^{-1}(z)$: each seed in that fiber has exactly one accompanying
$x=P_{m,r}^{-1}(y)$. This proves the bijectivity assertion.
Given a challenge $z=h_m(r)$, sample $Y\leftarrow U_m$ and run the
$F_m$-inverter on $(z,Y)$. Return its seed component after checking
$h_m(r')=z$. Conditional on $r$, $P_{m,r}(U_m)$ is uniform, so this is
exactly the joint distribution in an $F_m$ inversion experiment.
Whenever the full inversion succeeds, the recovered seed is correct.
In the reverse direction, recover $r$ from $z$ and then compute
$P_{m,r}^{-1}(y)$. Verify the seed before using it. Bijectivity of $h_m$
makes this successful exactly when the seed inverter is successful.
\hfill$\square$

This proves more than the failure of the unmasked example. In this
seed-indexed design, any hardness of the entire construction must
already reside in $h_m$. If $h_m$ has an efficient inverse, so does
$F_m$; if it does not, the missing hardness has merely moved to a
length-preserving bijection on the seed space. No average-case hardness
has been created by the trapdoor permutation in the second coordinate.
Polynomial comparability of the relevant length parameters would still
be required to translate this parameterized equivalence into the precise
all-length definition of the target.

\medskip\needspace{5\baselineskip}
\textbf{B. A conditional dense-domain construction with all lengths handled.}
Instead of exposing a generation tape, consider a deterministic,
recognizable domain of usable descriptions and data. The following
hypothesis makes the remaining work explicit.

\textbf{Hypothesis H (not established).}
There are fixed positive integers $r,s$, nonempty sets
$D_m\subseteq\{0,1\}^m$ for every $m\ge1$, and maps
$q_m:D_m\to D_m$ such that:
\begin{enumerate}[label=(\roman*),leftmargin=2.8em,itemsep=3pt]
\item Membership in $D_m$ and evaluation of $q_m$ on $D_m$ are performed
by fixed deterministic polynomial-time algorithms, and every $q_m$ is
bijective.
\item For all sufficiently large $m$, the density
$\rho_m=|D_m|/2^m$ satisfies $\rho_m\ge m^{-r}$.
\item For every probabilistic polynomial-time $C$, for all sufficiently
large $m$ (with the threshold allowed to depend on $C$),
\begin{equation}\label{eq:50:H}
 \Pr_{x\leftarrow\mathrm{Unif}(D_m),\,C}
       [C(1^m,q_m(x))=x]\le 1-m^{-s}.
\end{equation}
\end{enumerate}
The same $r,s$ work for every $C$. This hypothesis only asks for a weak
average-case failure rate on $D_m$, not negligible inversion success.
It does not contain a hidden seed or grant an oracle to the evaluator.
Sampling from $D_m$ is not an additional algorithmic assumption below;
its uniform distribution specifies the hardness experiment.

\textbf{Proposition B (conditional normalization theorem).}
Hypothesis~H implies a family meeting \emph{every} quantifier and domain
requirement in Problem~50. The construction is explicit relative to the
membership and evaluation algorithms in~H.

\textit{Step 1: totalize without losing bijectivity.}
Define
\begin{equation}\label{eq:50:guard}
 g_m(x)=\begin{cases}q_m(x),&x\in D_m,\\x,&x\notin D_m.\end{cases}
\end{equation}
Both pieces preserve their respective disjoint sets, so $g_m$ is a
uniform polynomial-time permutation of the whole cube. For any
inverter $C$, its success on a uniform input decomposes exactly into the
conditional successes on $D_m$ and its complement. Applying
\eqref{eq:50:H} to the restriction of $C$ gives, eventually,
\begin{equation}\label{eq:50:weak}
 \Pr[C(1^m,g_m(U_m))\ne U_m]
       \ge \rho_m m^{-s}\ge m^{-(r+s)}.
\end{equation}
The first inequality uses only failure on $D_m$; behavior outside it is
irrelevant. Nevertheless, simply returning the input inverts $g_m$ with
probability at least $1-\rho_m$. Thus totalization alone is not the
strong one-wayness demanded by the problem.

\textit{Step 2: a verified quantitative amplification lemma.}
The next argument is a self-contained direct-product proof, in the
classical weak-to-strong framework. \eref{50}{4}
It is important that it allows a \emph{joint} inverter on all blocks.

\textbf{Lemma C (random-position extraction).}
Let $g$ be an efficiently evaluable permutation of a finite binary cube
$\Omega$, let $k,R$ be positive integers, and let $A$ invert the
$k$-fold direct product on a uniform
challenge with probability $\epsilon$. If
$0<\epsilon_0\le\epsilon$, there is a repeated random-position inverter
$C$ for $g$ whose failure probability is at most
\begin{equation}\label{eq:50:extract}
 \frac{\log(2/\epsilon_0)}{k}
     +\exp\left(-\frac{R\epsilon_0}{2k}\right),
\end{equation}
using at most $R$ runs of $A$ and $O(Rk)$ evaluations of $g$.
For a power-of-two $k$, all random choices in this construction can be
sampled exactly using a bounded number of unbiased bits.

\textit{Proof.}
For a challenge tuple $\mathbf y\in\Omega^k$, let
\[
 a(\mathbf y)=\Pr[A(\mathbf y)=(g^{\times k})^{-1}(\mathbf y)]\in[0,1],
 \qquad \E a(\mathbf Y)=\epsilon.
\]
Here $(g^{\times k})^{-1}$ denotes the inverse of the direct-product map;
it is used to define success, not to compute it. For each $j$, put
\[
 b_j(y)=\E_{\mathbf Y_{-j}}a(Y_1,\ldots,Y_{j-1},y,Y_{j+1},\ldots,Y_k),
 \qquad b(y)=\frac1k\sum_{j=1}^k b_j(y).
\]
An extraction trial on $y$ chooses $j$ uniformly, samples independent
uniform \emph{inputs} for the other $k-1$ positions, evaluates $g$ there,
and inserts $y$ in position $j$. It runs $A$, checks every returned block
using $g$, and returns block $j$ only when \emph{all} checks pass.
Because $g$ is a permutation, the auxiliary challenges are exactly
uniform. This trial succeeds with probability exactly $b(y)$.

Set $\tau=\epsilon_0/(2k)$, let
$\mathcal B=\{y:b(y)<\tau\}$, and write
$\beta=\Pr[Y\in\mathcal B]$. No algorithm needs to recognize
$\mathcal B$. Splitting according to whether the tuple meets this bad
set and applying the union bound with the nonnegative weight $a$ gives
\begin{align}
 \epsilon
 &\le \Pr[\forall j,\ Y_j\notin\mathcal B]
       +\sum_{j=1}^k \E[a(\mathbf Y)1_{\{Y_j\in\mathcal B\}}]
       \label{eq:50:bad1}\\
 &= (1-\beta)^k+k\E[b(Y)1_{\{Y\in\mathcal B\}}]
 \le (1-\beta)^k+k\tau\beta
 \le (1-\beta)^k+\epsilon_0/2. \label{eq:50:bad2}
\end{align}
Independence is used only to evaluate the probability
that none of the independently sampled \emph{challenges} is bad.
There is no independence assertion about $A$'s answers.
Consequently
\[
 \epsilon_0/2\le(1-\beta)^k\le e^{-k\beta},
 \qquad \beta\le\log(2/\epsilon_0)/k.
\]
For any $y\notin\mathcal B$, $R$ trials with fresh independent coins
fail with probability at most
$(1-\tau)^R\le e^{-R\tau}$. Averaging over $y$ proves
\eqref{eq:50:extract}. Each trial is verifiable without knowing the
missing inverse, so the repeated algorithm stops only on a correct
answer. \hfill$\square$

\textit{Step 3: instantiate the parameters and the quantifiers.}
Set $t=r+s$ and define integer-valued functions
\begin{equation}\label{eq:50:lengths}
 k_m=2^{\lceil\log_2(m^{t+2})\rceil},\qquad
 M_m=m k_m,\qquad
 T_m(x_1,\ldots,x_{k_m})=(g_m(x_1),\ldots,g_m(x_{k_m})).
\end{equation}
The next power of two can be computed with integer arithmetic; no
floating-point logarithm is needed. We have
\[
 m^{t+2}\le k_m<2m^{t+2},\qquad
 m^{t+3}\le M_m<2m^{t+3}.
\]
The numbers $M_m$ are strictly increasing. The power-of-two choice for
$k_m$ makes position sampling use exactly $\log_2 k_m$ random bits,
avoiding an unbounded rejection-sampling subroutine.

Suppose a fixed uniform polynomial-time inverter $A$ for $T_m$ has
success at least $m^{-c}$ for infinitely many $m$, where $c$ is a fixed
positive integer. Apply Lemma~C with
\[
 \epsilon_0=m^{-c},\qquad R=2k_m m^{c+t+2}.
\]
At every such length its derived inverter has failure at most
\begin{equation}\label{eq:50:contradiction}
 \frac{\log 2+c\log m}{k_m}+e^{-m^{t+2}}
       < m^{-t}
\end{equation}
for all sufficiently large $m$. The algorithm does not estimate
$\epsilon$ or know which lengths are favorable. It uses the displayed
parameters at every length; its exponent $c$ is a fixed constant, not
advice depending on $m$. Its running time is polynomial in $m$, since
$M_m,k_m,R$ are. Infinitely many violations of \eqref{eq:50:weak}
contradict that bound's eventual validity for this one inverter.
Thus the success of every uniform polynomial-time inverter for $T_m$
is negligible in $m$.

\textit{Step 4: cover every binary length, without an easy subsequence.}
For every $N\ge1$, compute
\[
 m(N)=\max\{m\ge1:M_m\le N\}.
\]
This exists because $M_1=1$; it can be computed by integer search in
polynomial time. Split an $N$-bit string uniquely as
$u\Vert z$ with $|u|=M_{m(N)}$ and define
\begin{equation}\label{eq:50:alllengths}
 f_N(u\Vert z)=T_{m(N)}(u)\Vert z.
\end{equation}
This is a bijection on the \emph{entire} $N$-cube for every $N$, computed
by one polynomial-time algorithm. The unchanged padding is independent
uniform data, not a key-generation tape.

The security transfer needs a little care: one may not secretly give
an inverter a favorable length as advice. Suppose a fixed inverter $A$
for \eqref{eq:50:alllengths} succeeds at least $N^{-c}$ for infinitely
many $N$. On a challenge for $T_m$, a single uniform inverter runs $A$
for \emph{every} integer
\[
 N\in I_m:=\{M_m,M_m+1,\ldots,M_{m+1}-1\},
\]
using a fresh uniform padding string of length $N-M_m$ in each run.
It verifies the proposed prefix inverse under $T_m$ and returns the
first verified one. This loop is polynomial in $m$:
$|I_m|\le M_{m+1}<2(m+1)^{t+3}$, and each run has polynomial input length.
For each fixed $N\in I_m$, the augmented challenge has exactly the
uniform $N$-bit distribution. The loop therefore succeeds with at least
$A$'s success at that $N$. Infinitely many favorable $N$ give infinitely
many different $m$, with success at least
\[
 \bigl(2(m+1)^{t+3}\bigr)^{-c},
\]
a nonnegligible function of $m$. This contradicts Step~3. No favorable
length has been selected nonuniformly, and both the target's
all-length requirement and its universal adversary quantifier are
satisfied. This completes the proof of Proposition~B, \emph{conditional
on H}. \hfill$\square$

\medskip\needspace{5\baselineskip}
\textbf{C. Trying to instantiate H, and finding its exact failure points.}
A tempting choice is to let $D_m$ consist of valid public descriptions
and payloads and let $q_m$ preserve the public description while
permuting the payload. For this to instantiate H one must prove, in a
single fixed encoding, polynomial-time membership in $D_m$, inverse-
polynomial density among all $m$-bit strings, invariance of $D_m$ under
the proposed map, and the bound \eqref{eq:50:H} for the uniform
conditional distribution. Security under a key generator's distribution
does not by itself imply security under uniform valid encodings.
Repeated encodings and highly unequal sampling weights can change that
distribution. Adding a generation witness makes membership easier, but
when that witness is the complete seed, Lemma~A gives an inverse instead
of hardness. Thus the construction does not currently supply an
instance of H from the cited keyed result.

Inverse-polynomial density is not a cosmetic assumption. If $\rho_m$
is negligible, any polynomially bounded number $k(m)$ of copies of the
guarded map \eqref{eq:50:guard} has an identity-inverter success of at
least
\begin{equation}\label{eq:50:sparse}
 (1-\rho_m)^{k(m)}\ge1-k(m)\rho_m.
\end{equation}
The proof is that all blocks outside $D_m$ are fixed, independently,
and the final inequality is the union bound. For example,
$\rho_m=2^{-m/2}$ cannot be amplified by polynomially many copies in
this construction, no matter how difficult inversion on $D_m$ might be.
An infinite search for a single difficult input, or a sparse collection
of difficult encodings, would not repair the missing density condition.

There is also a concrete fatal error in a common simplified product
proof. Even for the identity permutation, an algorithm may flip one
coin and either output the entire correct tuple or deliberately fail
every block. The joint success is $1/2$, not $2^{-k}$, although each
individual block is correct with probability $1/2$. Lemma~C was derived
to remove exactly this unjustified independence step.

\medskip\needspace{5\baselineskip}
\textbf{D. Structural construction attempts: gates, order, and cycles.}
A natural way to turn an arbitrary difficult function into a permutation
is a Feistel round. For any efficiently computable
$h:\{0,1\}^d\to\{0,1\}^d$, define
\[
 \mathcal F_h(L,R)=(R,L\mathbin\oplus h(R)).
\]
But
\begin{equation}\label{eq:50:feistel}
 \mathcal F_h^{-1}(U,V)=(V\mathbin\oplus h(U),U).
\end{equation}
Only a \emph{forward} evaluation of $h$ is used by the inverse. A
publicly specified polynomial sequence of such rounds is inverted by
reversing the sequence, even if the round functions themselves were
one-way. This attempt guarantees bijectivity too transparently.

The same defect appears in a putative universal clean reversible
implementation. If a uniformly constructible polynomial-size circuit
$R_n$ over a fixed set of explicitly invertible gates satisfies
\[
 R_n(x,0^{a(n)})=(f_n(x),0^{a(n)}),
\]
then reversing its gates on $(y,0^{a(n)})$ computes $f_n^{-1}(y)$ in
polynomial time. Ordinary forward computation of $f_n$ does not justify
this \emph{clean} implementation: retaining $(x,f_n(x))$ leaves the very
preimage that the inverse problem asks to recover. No claim about all
polynomial-time bijections follows from the conditional circuit argument.

An order-preserving construction is also too rigid. Give the Boolean
cube its coordinatewise partial order. Every order-preserving bijection
of this finite poset has an order-preserving inverse: some finite power
of the bijection is the identity, so the inverse is a positive power.
Hence it is an order automorphism, permutes the atoms $e_1,\ldots,e_n$,
and preserves joins. Since every $x$ is the join of its supported atoms,
\[
 f(x)=\bigvee_{i:x_i=1}f(e_i)
\]
is merely a coordinate permutation. The inverse can be learned with
$n$ forward evaluations and then applied efficiently. The finite order
used in the proof need not be found by the inversion algorithm.

Finally, let $\ell_f(y)$ be the length of the cycle containing $y$.
Starting with $y$, apply $f$ until the first return to $y$. If the cycle
length is at most a given cutoff $T$, the state immediately before that
return is $f^{-1}(y)$. Consequently the target requires, for every
polynomial $T(n)$,
\begin{equation}\label{eq:50:cycles}
 \Pr_{y\leftarrow U_n}[\ell_{f_n}(y)\le T(n)]
        \quad\text{to be negligible.}
\end{equation}
This condition is necessary, not sufficient: binary addition
$x\mapsto x+1\pmod{2^n}$ is one cycle of length $2^n$, yet subtraction
inverts it in linear time. Efficiently invertible conjugation cannot
help. If $\widetilde f=h^{-1}\circ f\circ h$ with uniformly efficient
$h,h^{-1}$, then the two inversion experiments reduce to each other
with equal success, because $h(U_n)$ is uniform; their cycle lengths
also agree. Long cycles and relabeling have not supplied a hard domain.

\medskip\needspace{5\baselineskip}
\textbf{E. A dependency check using unique inverse bits.}
For any full-cube family $f$ as in the target's evaluation requirement,
define the total binary language, using any fixed unambiguous tuple
encoding,
\[
 L_f=\{\langle1^n,y,i\rangle: |y|=n,\ 1\le i\le n,
                         \ (f_n^{-1}(y))_i=1\}.
\]
Malformed encodings are outside $L_f$. Recall that $\mathsf{UP}$ consists
of languages having a polynomial-time nondeterministic decider with at
most one accepting computation on each input; $\mathsf{coUP}$ consists
of their complements.
On a well-formed tuple, guess exactly $n$ bits $x$, verify $f_n(x)=y$,
and check $x_i=1$. There is at most one accepting witness. For the
complement, use $x_i=0$ on well-formed tuples and one deterministic
accepting branch on malformed ones. Thus
\begin{equation}\label{eq:50:up}
 L_f\in\mathsf{UP}\cap\mathsf{coUP}.
\end{equation}
If $L_f$ were in $\mathsf P$, $n$ bit queries would invert $f_n$ on
\emph{every} output. If $L_f$ were in $\mathsf{BPP}$, amplify each bit
query's error to at most $1/(10n)$ and take a union bound; this yields a
uniform randomized inverter with success at least $9/10$ on every
output. Therefore a family meeting Problem~50 implies
\[
 \mathsf{UP}\cap\mathsf{coUP}\not\subseteq\mathsf{BPP},
 \qquad
 \mathsf P\ne\mathsf{NP}\cap\mathsf{coNP},
 \qquad \mathsf P\ne\mathsf{NP}.
\]
The latter two consequences follow from
$\mathsf{UP}\cap\mathsf{coUP}\subseteq\mathsf{NP}\cap\mathsf{coNP}$
and $\mathsf P\subseteq\mathsf{BPP}$. Existence also implies existence
of ordinary one-way functions simply by forgetting bijectivity.
Thus an unconditional instantiation of H would settle more than a domain
encoding problem: it would also settle catalogue Problems~1, 8, and~34.
These implications are proved here; no reverse implication from a
worst-case separation is being used.

A finer error audit illustrates the missing average-case estimate.
Suppose a fixed randomized predictor $B(1^n,y,i)$ has average bit error
\[
 e_n=\frac1n\sum_{i=1}^n
      \Pr_{y\leftarrow U_n,B}[B(1^n,y,i)\ne(f_n^{-1}(y))_i].
\]
Running it once for each bit gives a complete candidate whose failure
probability is at most $n e_n$, regardless of correlations among bit
errors. Hence the target forces
$e_n\ge(1-\eta(n))/n$ for a negligible $\eta$ depending on this
predictor. This is only a necessary estimate; small advantage over
random bit guessing does not automatically recover the full inverse.
Neither \eqref{eq:50:up} nor this estimate proves an average-case lower
bound for any explicit $f$.

\paragraph{Stress test.}
\textit{Exact target and source fit.}
A sampled-key theorem has not been relabeled as a key-free one. Keeping
a generation tape and masking it were separately tested; their inverses
and reductions were supplied. The indexed versus single-function issue
is independently documented in Goldreich's clarification.
\eref{50}{1} The present conditional construction assumes H outright;
it does not claim a fully black-box derivation of H from arbitrary
one-way functions or from obfuscation, and therefore does not bypass
Asharov--Segev's model-specific separation. \eref{50}{2}

\textit{All lengths, randomness, and adversaries.}
Every finite-length map in \eqref{eq:50:alllengths} is exactly bijective,
including at small lengths where the hardness assumption says nothing.
The padding proof enumerates polynomially many lengths instead of
hardwiring favorable ones. The failure bounds hold eventually for each
fixed derived inverter, which is why an infinite set of good attack
lengths is enough for a contradiction. Position sampling is exact with
bounded randomness, and every trial's output is checked. All
expectations and sums in Lemma~C are finite; no limit exchange is used.
The only independence assumptions are explicitly generated independent
challenge blocks or fresh repeated trials, never independent answers
from a joint adversary.

\textit{Finite falsification tests.}
The companion script \texttt{50\_checks.py} performs exact-arithmetic
checks, not security experiments. It checks every function on a
four-element set in the Feistel identity, every subset and permutation
of that subset for the guarded extension on a four-element universe,
and every three-block $\{0,1\}$ success profile with values in
$\{0,\tfrac12,1\}$ for the weighted bad-set inequality in Lemma~C.
It also exhausts all permutations of the three-dimensional Boolean
cube for the order-preserving classification, tests the masked-seed
fiber-count identity including noninjective masks, and checks the
all-length block partition on a finite range. Counts and outcomes are
recorded in \texttt{50\_checks.json}. These tests can expose small
counterexamples or implementation errors; they are not evidence that
any explicit family is hard against every polynomial-time adversary.
The displayed proofs, not the tested cutoffs, justify the general
lemmas.

\textit{What remains missing.}
No explicit pair $(D_m,q_m)$ has been shown to satisfy H's hardness
condition. In particular, H must not be presented as an independently
easy conjecture: allowing $D_m$ to be the whole cube shows that the
existence of such a primitive is also implied by the target itself.
Proposition~B is a normalization equivalence at the level of existence,
with useful explicit domain and security bookkeeping, not a new
unconditional lower bound. The order, Feistel, cycle, and seed
constructions fail for the specific reasons proved above; they do not
exhaust all polynomial-time bijections. No counterexample to the
existence statement has been constructed.

\paragraph{Outcome.}
\textbf{Outcome: Conditional reduction.}

An explicit guarded-extension, direct-product, and all-length-padding
construction was proved conditional on a uniformly recognizable dense
partial permutation with a fixed inverse-polynomial inversion failure
rate. The quantitative extraction argument handles joint attacks and
the padding argument preserves uniformity. The attempt also proves an
exact masked-seed security equivalence and records several constructive
inversion obstructions. These are auxiliary deductions and a
self-contained use of classical amplification, not asserted historical
novelties.

A complete positive resolution would still require an unconditional
construction satisfying H, or a substantially different explicit
full-cube permutation with a proof of negligible inversion success.
A negative resolution would require ruling out \emph{all} families in
the original statement, not just the tested construction classes. Neither
has been obtained here.
% END RESEARCH INSERTION 50
\subsection{Sources}
\begin{itemize}\raggedright
\item[\textnormal{[S1]}]\hypertarget{source:50:1}{} On Constructing One-Way Permutations from Indistinguishability Obfuscation, ePrint2015/752, TCC2016 version.
\url{https://eprint.iacr.org/2015/752.pdf}.
{\small\textit{Catalogue formulation source.}}
\item[\textnormal{[S2]}]\hypertarget{source:50:2}{} On One-Shot Signatures, Quantum vs Classical Binding, and Obfuscating Permutations.
\url{https://arxiv.org/abs/2507.12456}.
{\small\textit{Catalogue research/status reference; not a proof endorsement.}}
\item[\textnormal{[S3]}]\hypertarget{source:50:3}{} Quantum Search with In-Place Queries.
\url{https://drops.dagstuhl.de/storage/00lipics/lipics-vol350-tqc2025/LIPIcs.TQC.2025.1/LIPIcs.TQC.2025.1.pdf}.
{\small\textit{Catalogue research/status reference; not a proof endorsement.}}
% BEGIN EXTRA SOURCES 50
\item[\textnormal{[E1]}]\hypertarget{extra:50:1}{} Oded Goldreich,
\emph{On one-way permutations (a clarification)}, November 2005,
Cases 1--3, especially Case 2 and the distinction between indices and
index-generation coins.
\url{https://www.wisdom.weizmann.ac.il/~oded/owp.html}.
{\small\textit{Author's primary clarification; consulted 27 September 2026.}}

\item[\textnormal{[E2]}]\hypertarget{extra:50:2}{} Gilad Asharov and Gil Segev,
\emph{On Constructing One-Way Permutations from Indistinguishability
Obfuscation}, Journal of Cryptology \textbf{31} (2018), 698--736,
DOI: 10.1007/s00145-017-9268-6; earlier ePrint 2015/752 and TCC 2016-A.
\url{https://link.springer.com/article/10.1007/s00145-017-9268-6}.
Author presentation, slides 28 and 44 (one-based page numbers):
\url{https://www.cs.cornell.edu/~asharov/slides/AS16.pdf}.
{\small\textit{Publisher abstract and author slides consulted 27 September
2026 for the fully black-box/domain-invariance restrictions. Direct
retrieval of the ePrint PDF failed; a full-paper proof audit is not
claimed. This entry supplements, rather than replaces, [S1].}}

\item[\textnormal{[E3]}]\hypertarget{extra:50:3}{} Omri Shmueli and Mark Zhandry,
\emph{On One-Shot Signatures, Quantum vs Classical Binding, and Obfuscating
Permutations}, arXiv:2507.12456v1, 16 July 2025. Theorem 6; Section 1.4;
Section 5.1, Construction 41 and Theorem 42.
\url{https://arxiv.org/html/2507.12456v1}.
{\small\textit{Specified manuscript version and relevant construction/proof
passages consulted 27 September 2026. Conditional result; neither its
assumptions nor an independent verification of the whole paper is
asserted here.}}

\item[\textnormal{[E4]}]\hypertarget{extra:50:4}{} Oded Goldreich,
\emph{Foundations of Cryptography, Volume 1: Basic Tools},
Cambridge University Press, 2001, Section 2.3,
``Weak One-Way Functions Imply Strong Ones,'' especially Section 2.3.1,
``The Construction and its Analysis.'' Author's detailed section guide:
\url{https://www.wisdom.weizmann.ac.il/~oded/foc-vol1.html}.
{\small\textit{Author's section guide consulted 27 September 2026 to
attribute the classical amplification principle. The finite-instance
bound used in this notebook is proved in full above, not imported as
an unchecked quotation or theorem-number claim.}}

\item[\textnormal{[E5]}]\hypertarget{extra:50:5}{} Blake Holman, Ronak
Ramachandran, and Justin Yirka, \emph{Quantum Search with In-Place Queries},
in 20th Conference on the Theory of Quantum Computation, Communication
and Cryptography (TQC 2025), LIPIcs \textbf{350} (2025), 1:1--1:18,
Theorem 2 and Section 3. DOI: 10.4230/LIPIcs.TQC.2025.1.
\url{https://drops.dagstuhl.de/storage/00lipics/lipics-vol350-tqc2025/html/LIPIcs.TQC.2025.1/LIPIcs.TQC.2025.1.html}.
{\small\textit{Published HTML text consulted 27 September 2026. A quantum
query-model result, not an unconditional classical one-way-permutation
construction. This entry supplies authors and the exact result for [S3].}}
% END EXTRA SOURCES 50
\end{itemize}

% END ORIGINAL SECTION 50
\end{document}
